On considère la fonction $f$ définie sur $\R$ par
\[ f(x) = x^2+2x-3. \]

La courbe représentative de la fonction $f$ dans un repère orthogonal se note $\mathcal{P}$.

\begin{enumerate}
	\item La courbe $\mathcal{P}$ est une parabole. Déterminer les coordonnées de son sommet.

	\item On considère la fonction $g$ définie sur $\R$ par
	\[ g(x) = x^2+x-6. \]
	\begin{enumerate}
		\item Justifier que $-3$ et $2$ sont racines de la fonction $g$.
		\item Déterminer le tableau de signes de la fonction $g$ sur $\R$.
	\end{enumerate}

	\item On considère la fonction $h$ définie sur $\R$ par
	\[ h(x) = x+3. \]
	La courbe représentative de la fonction $h$ se note $\mathcal{D}$.

	Vérifier que pour tout réel $x$, on a $f(x) - h(x) = g(x)$.

	\item En déduire la position relative des courbes $\mathcal{P}$ et $\mathcal{D}$.

	On déterminera les coordonnées des éventuels points d'intersection.

	\item Réaliser un schéma donnant l'allure des courbes $\mathcal{P}$ et $\mathcal{D}$ sur lequel figureront les résultats établis à la question \textbf{1.} et à la question \textbf{4.}
\end{enumerate}
